\documentclass[mnsc,nonblindrev]{informs3_NoOR} 
\OneAndAHalfSpacedXI
\usepackage{thmtools}
\usepackage[dvipsnames]{xcolor}
\usepackage{verbatim}
\usepackage{multirow}
\usepackage{float} 
\usepackage{endnotes}   
\usepackage{enumitem}   	 
\usepackage{amsmath}
\usepackage{mathrsfs}
\usepackage{mathtools}
\usepackage{url} 
\usepackage{subcaption} 
\usepackage[T1]{fontenc}
\usepackage[english]{babel}
\usepackage{lmodern} 
\usepackage{bbm}
\usepackage{lscape}
\usepackage{scalerel}
\usepackage{xr, xr-hyper}
\usepackage{hyperref}
\usepackage[colorinlistoftodos,prependcaption,textsize=tiny]{todonotes}
\usepackage{regexpatch}
\makeatletter
\xpatchcmd{\@todo}{\setkeys{todonotes}{#1}}{\setkeys{todonotes}{inline,#1}}{}{}
\makeatother
\hypersetup{
    colorlinks,
    linkcolor={red!50!black},
    citecolor={blue!50!black},
    urlcolor={blue!80!black}
}
\usepackage{natbib}
 \bibpunct[, ]{(}{)}{,}{a}{}{,}%
 \def\bibfont{\small}%
 \def\bibsep{\smallskipamount}%
 \def\bibhang{24pt}%
 \def\newblock{\ }%
 \def\BIBand{and}%
\TheoremsNumberedThrough 
\ECRepeatTheorems
\EquationsNumberedThrough   
\MANUSCRIPTNO{NoNumber}
\usepackage{soul}
\sethlcolor{yellow}

\begin{document}

\TITLE{Content Generation: Human Creators, AI Learning, and Platform Design}
\maketitle


\section{Introduction}

\section{Related Literature}
\begin{itemize}
\item \cite{ai_genai_2026}
\item \cite{esmaeili_how_2025}
\item \cite{taitler_braesss_2024}
\item \cite{taitler_data_2025}
\item \cite{taitler_selective_2025}
\item \cite{yao_rethinking_2023}
\item \cite{yao_human_2024}
\item \cite{yao_user_2024}
\item \cite{yu_beyond_2025}
\end{itemize}


%%%%%%%%%%%%%%%%%%%%%%%%%%%%%%%%%%%%%%%%%%%%%%%%%%%%%%%%%%
\section{Model}\label{sec:model}
%%%%%%%%%%%%%%%%%%%%%%%%%%%%%%%%%%%%%%%%%%%%%%%%%%%%%%%%%%
 
We analyze a platform model in which a human creator and an AI content generator compete for consumer attention. The platform sets strategic parameters in Stage~1; the human creator chooses effort in Stage~2. We solve by backward induction.
 
\subsection{Content Generation and Quality}
 
The human creator exerts effort $q_h \geq 0$ at cost $C(q_h) = \tfrac{1}{2}q_h^2$. The AI generates content of quality
\begin{equation}\label{eq:qA}
q_A = \alpha\, q_h + Q,
\end{equation}
where $\alpha \in (0,1]$ is the AI's learning efficiency and $Q \geq 0$ is its baseline quality.
 
\subsection{Consumer Demand}
 
Consumers are uniformly distributed on $[0,1]$, with location~$x$ indexing proximity to the AI (at~0) relative to the human (at~1). Consumer utility is
\begin{equation}\label{eq:utility}
u_A = q_A - t\,x\,m_A, \qquad u_h = q_h - t(1-x)\,m_H,
\end{equation}
where $t > 0$ is the mismatch cost, $\beta_h \in [0,1]$ is the platform's promotion weight for human content, $\delta \in (0,1)$ scales algorithmic influence, and
\begin{equation}\label{eq:mismatch}
m_A(\beta_h) = 1 - \delta(1-\beta_h), \qquad m_H(\beta_h) = 1 - \beta_h\delta.
\end{equation}
Both lie in $[1-\delta,\,1]$. Increasing $\beta_h$ raises $m_A$ (penalizing AI) and lowers $m_H$ (favoring human). The demand functions are
\begin{equation}\label{eq:demand}
D_A = \frac{q_A}{t\,m_A}, \qquad D_h = \frac{q_h }{t\,m_H}.
\end{equation}
 
\subsection{Platform and Creator Payoffs}
 
The platform incurs a discrimination penalty $\tfrac{k}{2}(\beta_h - \tfrac{1}{2})^2$, $k \geq 0$, representing the regulatory or reputational cost of departing from neutral curation. Under view-based revenue, the platform's payoff is
\begin{equation}\label{eq:up-view}
u_p^V = D_A + D_h -r\,q_h - \tfrac{k}{2}(\beta_h - \tfrac{1}{2})^2.
\end{equation}
Under engagement-based revenue it is
\begin{equation}\label{eq:up-eng}
u_p^E = q_A D_A + q_h D_h - r q_h - \tfrac{k}{2}(\beta_h - \tfrac{1}{2})^2.
\end{equation}
The human creator earns
\begin{equation}\label{eq:uc}
u_c = r\,q_h - \tfrac{1}{2}q_h^2.
\end{equation}

%%%%%%%%%%%%%%%%%%%%%%%%%%%%%%%%%%%%%%%%%%%%%%%%%%%%%%%%%%
\section{Equilibrium Analysis: View-Based Revenue}\label{sec:view}
%%%%%%%%%%%%%%%%%%%%%%%%%%%%%%%%%%%%%%%%%%%%%%%%%%%%%%%%%%
 
\subsection{Stage 2: Creator's Optimal Effort}
 
\begin{proposition}\label{prop:effort}
For any compensation rate $r \geq 0$, the creator's optimal effort is $q_h^* = r$, yielding AI quality $q_A^* = \alpha r + Q$. The creator's equilibrium welfare is $u_c^* = \tfrac{1}{2}r^2 \geq 0$, so participation is always individually rational.
\end{proposition}
 
\subsection{Stage 1: Optimal Compensation}
 
Define $W(\beta_h) = \frac{\alpha}{tm_A} + \frac{1}{tm_H}$, the composite demand sensitivity.
 
\begin{proposition}\label{prop:rstar-view}
The platform's payoff is strictly concave in~$r$ (the second derivative is $-2$), with unique maximiser
\begin{equation}\label{eq:rstar-view}
r^* = \frac{W}{2} = \frac{1}{2}\!\left(\frac{\alpha}{tm_A} + \frac{1}{tm_H}\right).
\end{equation}
The equilibrium qualities are $q_h^* = W/2$ and $q_A^* = \alpha W/2 + Q$.
\end{proposition}
 
\subsection{Stage 1: Optimal Promotion}
 
Let $y_A = 1/m_A$, $y_H = 1/m_H$, and define the marginal revenue from promotion as
\begin{equation}\label{eq:L-def}
L(\beta_h) = \frac{\delta}{t}\!\left[\frac{r^*}{m_H^2} - \frac{q_A^*}{m_A^2}\right] = \frac{\delta}{t}\!\left[\frac{W/2}{m_H^2} - \frac{\alpha W/2 + Q}{m_A^2}\right].
\end{equation}
Let $\Gamma(\beta_h) = \frac{1}{2}(W')^2 + \frac{1}{2}WW'' + \frac{Q}{t}y_A''$ be the second derivative of the unpenalised utility function.
 
\begin{proposition}\label{prop:beta-view}
The penalised utility function is
\begin{equation}\label{eq:VV}
U^V(\beta_h) = \frac{1}{4}W^2 + \frac{Q}{t}\,y_A - \frac{k}{2}\!\left(\beta_h - \frac{1}{2}\right)^{\!2}.
\end{equation}
Since $\Gamma(\beta_h) > 0$ for all $\beta_h \in [0,1]$, the unpenalised utility function is always strictly convex.  Let $\bar{\Gamma} := \max_{\beta_h \in [0,1]} \Gamma(\beta_h)$.
\begin{enumerate}
\item[\textup{(i)}] \textbf{$k = 0$:} The optimum is always a corner: $\beta_h^* \in \{0,1\}$, chosen by comparing $U^V(0)$ and $U^V(1)$. There exists a critical threshold $Q^*$ such that $\beta_h^* = 1$ for $Q < Q^*$ and $\beta_h^* = 0$ for $Q > Q^*$.
 
\item[\textup{(ii)}] \textbf{$k > \bar\Gamma$:} $U^V$ is strictly concave and the maximum is unique. Define $k^* = \max\!\big\{\bar\Gamma,\; -2L(0),\; 2L(1)\big\}$.
\begin{enumerate}
\item[\textup{(a)}] If $k > k^*$: the unique maximum is interior, $\beta_h^* \in (0,1)$, satisfying $L(\beta_h^*) = k(\beta_h^* - \frac{1}{2})$.
\item[\textup{(b)}] If $\bar\Gamma < k \leq k^*$ and $L(0) \leq -k/2$: $\beta_h^* = 0$.
\item[\textup{(c)}] If $\bar\Gamma < k \leq k^*$ and $L(1) \geq k/2$: $\beta_h^* = 1$.
\end{enumerate}
 
\item[\textup{(iii)}] \textbf{$0 < k \leq \bar\Gamma$:} $U^V$ may change concavity. The global maximum is found by comparing all stationary points of $L(\beta_h) = k(\beta_h - \frac{1}{2})$ with $U^V(0)$ and $U^V(1)$.
\end{enumerate}
\end{proposition}
 
%%%%%%%%%%%%%%%%%%%%%%%%%%%%%%%%%%%%%%%%%%%%%%%%%%%%%%%%%%
\section{Equilibrium Analysis: Engagement-Based Revenue}\label{sec:eng}
%%%%%%%%%%%%%%%%%%%%%%%%%%%%%%%%%%%%%%%%%%%%%%%%%%%%%%%%%%
 
\subsection{Stage 2: Creator's Optimal Effort}
 
The creator's problem is unchanged; Proposition~\ref{prop:effort} applies: $q_h^* = r$ and $q_A^* = \alpha r + Q$.
 
\subsection{Stage 1: Optimal Compensation}
 
Define $\Delta(\beta_h) = tm_Am_H - \alpha^2 m_H - m_A$.
 
\begin{proposition}\label{prop:rstar-eng}
The second derivative of $u_p^E$ with respect to $r$ is $2(\alpha^2/(tm_A) + 1/(tm_H) - 1)$, with sign determined by $\Delta$:
\begin{enumerate}
\item[\textup{(a)}] \textbf{$\Delta > 0$ (concavity):} The unique maximiser is
\begin{equation}\label{eq:rstar-eng}
r^* = \frac{\alpha\,Q\,m_H}{\Delta}.
\end{equation}
$r^* > 0$ if and only if $Q > 0$; when $Q = 0$, $r^* = 0$ and the creator is inactive. A sufficient condition for $\Delta > 0$ at all $\beta_h$ is $t > \alpha^2 + \frac{1}{1-\delta}$.
 
\item[\textup{(b)}] \textbf{$\Delta \leq 0$ (convexity):} The payoff is non-concave in $r$ and the platform's optimal subsidy is unbounded. The engagement formulation requires $\Delta > 0$ for well-defined solutions.
\end{enumerate}
\end{proposition}
 
\subsection{Stage 1: Optimal Promotion}
 
Assume $\Delta(\beta_h) = tm_Am_H - \alpha^2 m_H - m_A > 0$ and $Q > 0$ throughout, so the engagement payoff is concave in $r$ and $r^* = \alpha Qm_H/\Delta > 0$. Define the marginal promotion effect
\begin{equation}\label{eq:F-def}
L(\beta_h) = \frac{\delta}{t}\!\left[\frac{(r^*)^2}{m_H^2} - \frac{(q_A^*)^2}{m_A^2}\right] = \frac{\delta\,Q^2}{t\,\Delta^2}\!\left[\alpha^2 - (tm_H - 1)^2\right],
\end{equation}
and the promotion crossover point $\beta_c = (t - 1 - \alpha)/(t\delta)$. $F > 0$ iff $\beta_h > \beta_c$.
 
\begin{proposition}[Optimal Promotion: Engagement-Based]\label{prop:beta-eng}
Assume $\Delta > 0$ and $Q > 0$. The penalised value function is
\begin{equation}
    U^E(\beta_h) = \frac{Q^2(tm_H - 1)}{t\,\Delta} - \frac{k}{2}\!\left(\beta_h - \frac{1}{2}\right)^{\!2}. \label{eq:VE}
\end{equation}


The marginal promotion effect $L(\beta_h) = \frac{\delta Q^2}{t\Delta^2}[\alpha^2 - (tm_H-1)^2]$ satisfies $L > 0$ iff $\beta_h > \beta_c$, where $\beta_c = (t-1-\alpha)/(t\delta)$. The optimal promotion depends on the mismatch cost $t$ and the penalty $k$:
 
\begin{enumerate}
\item[\textup{(i)}] \textbf{$k = 0$.} The optimum is always a corner: $\beta_h^* \in \{0,1\}$.
\begin{enumerate}[label=(\alph*)]
\item If $t \leq 1+\alpha$: $\beta_h^* = 1$.
\item If $t \geq (1+\alpha)/(1-\delta)$: $\beta_h^* = 0$.
\item If $1+\alpha < t < (1+\alpha)/(1-\delta)$: $\beta_h^* \in \{0,1\}$, chosen by comparing $U^E(0)$ and $U^E(1)$.
\end{enumerate}
 
\item[\textup{(ii)}] \textbf{$k > 0$, $t \leq 1+\alpha$ (positive effect dominates).} Let $k_A = 2L(1)$.
\begin{enumerate}[label=(\alph*)]
\item If $k \leq k_A$: $\beta_h^* = 1$.
\item If $k > k_A$: $\beta_h^* \in (\frac{1}{2}, 1)$, satisfying $F(\beta_h^*) = k(\beta_h^* - \frac{1}{2})$.
\end{enumerate}
 
\item[\textup{(iii)}] \textbf{$k > 0$, $t \geq (1+\alpha)/(1-\delta)$ (negative effect dominates).} Let $k_B = 2|L(0)|$.
\begin{enumerate}[label=(\alph*)]
\item If $k \leq k_B$: $\beta_h^* = 0$.
\item If $k > k_B$: $\beta_h^* \in (0, \frac{1}{2})$, satisfying $F(\beta_h^*) = k(\beta_h^* - \frac{1}{2})$.
\end{enumerate}
 
\item[\textup{(iv)}] \textbf{$k > 0$, $1+\alpha < t < (1+\alpha)/(1-\delta)$ (sign-changing regime).} Here $L(0) < 0$ and $L(1) > 0$. Let $k^{\star} = \max\{2|L(0)|,\, 2L(1)\}$.
\begin{enumerate}[label=(\alph*)]
\item If $k > k^{\star}$: the global maximum is interior, $\beta_h^* \in (0,1)$.
\item If $k \leq \min\{2|L(0)|,\, 2L(1)\}$:  $\beta_h^*$ is either a corner or at an interior stationary point.
\item If $k \leq 2|L(0)|$ and $k > 2L(1)$: $\beta_h^*$ is either $0$ or at an interior stationary point.
\item If $k > 2|L(0)|$ and $k \leq 2L(1)$: $\beta_h^*$ is either $1$ or at an interior stationary point.
\end{enumerate}
\end{enumerate}
 
In all cases, as $k \to \infty$, $\beta_h^* \to \frac{1}{2}$.
 
When $Q = 0$: $r^* = 0$, $U^E = -\frac{k}{2}(\beta_h - \frac{1}{2})^2$, and $\beta_h^* = \frac{1}{2}$.
\end{proposition} 
 
\bibliographystyle{plainnat}
\bibliography{references}
 
%%%%%%%%%%%%%%%%%%%%%%%%%%%%%%%%%%%%%%%%%%%%%%%%%%%%%%%%%%
\appendix
\section{Proofs}\label{app:proofs}
%%%%%%%%%%%%%%%%%%%%%%%%%%%%%%%%%%%%%%%%%%%%%%%%%%%%%%%%%%
 
\subsection*{Proof of Proposition~\ref{prop:effort}}
 
The creator maximises $u_c = rq_h - \frac{1}{2}q_h^2$. The first-order condition is $\partial u_c/\partial q_h = r - q_h = 0$, giving $q_h^* = r$. The second-order condition $\partial^2 u_c/\partial q_h^2 = -1 < 0$ confirms a global maximum for any $r \geq 0$. Substituting into~\eqref{eq:qA}: $q_A^* = \alpha r + Q$. Creator welfare: $u_c^* = r \cdot r - \frac{1}{2}r^2 = \frac{1}{2}r^2 \geq 0$. \hfill$\square$
 
\subsection*{Proof of Proposition~\ref{prop:rstar-view}}
 
Substitute $q_h^* = r$ and $q_A^* = \alpha r + Q$ into the view-based payoff (omitting the penalty, which is $r$-independent):
\[
u_p^V = \frac{\alpha r + Q}{tm_A} + \frac{r}{tm_H} - r^2 = -r^2 + r\!\left(\frac{\alpha}{tm_A} + \frac{1}{tm_H}\right) + \frac{Q}{tm_A}.
\]
This is a quadratic $-r^2 + Wr + Q/(tm_A)$ with second derivative $-2 < 0$, confirming strict concavity for all parameter values. The unique maximiser is $r^* = W/2$. Since $W = \alpha/(tm_A) + 1/(tm_H) > 0$, $r^* > 0$ is automatic. \hfill$\square$
 
\subsection*{Proof of Proposition~\ref{prop:beta-view}}
 
\noindent\textbf{utility function.}
At $r^* = W/2$: $u_p^V = -(W/2)^2 + W \cdot W/2 + Q/(tm_A) = W^2/4 + Q/(tm_A)$. Adding the penalty:
\[
U^V(\beta_h) = \frac{1}{4}W^2 + \frac{Q}{t}\,y_A - \frac{k}{2}\!\left(\beta_h - \frac{1}{2}\right)^{\!2}.
\]
 
\medskip\noindent\textbf{Strict convexity of the unpenalised utility function.}
The derivatives of $y_A = 1/m_A$ and $y_H = 1/m_H$ are:
\[
y_A' = \frac{-\delta}{m_A^2},\quad y_A'' = \frac{2\delta^2}{m_A^3} > 0; \qquad y_H' = \frac{\delta}{m_H^2},\quad y_H'' = \frac{2\delta^2}{m_H^3} > 0.
\]
Since $m_A' = \delta > 0$, $y_A$ is decreasing; since $m_H' = -\delta < 0$, $y_H$ is increasing. The composite sensitivity is $W = \frac{\alpha}{t}y_A + \frac{1}{t}y_H$, with
\[
W' = \frac{\delta}{t}\!\left(\frac{1}{m_H^2} - \frac{\alpha}{m_A^2}\right), \qquad W'' = \frac{\alpha}{t}y_A'' + \frac{1}{t}y_H'' > 0.
\]
The unpenalised second derivative is
\[
\Gamma(\beta_h) = \frac{1}{2}(W')^2 + \frac{1}{2}WW'' + \frac{Q}{t}y_A''.
\]
Each term is non-negative: $(W')^2 \geq 0$; $WW'' > 0$ since $W > 0$ and $W'' > 0$; $\frac{Q}{t}y_A'' \geq 0$ since $Q \geq 0$. The middle term is strictly positive, so $\Gamma > 0$ for all $\beta_h$, all parameters.


 
\medskip\noindent\textbf{Part~(i): $k = 0$.}
$V''^V = \Gamma > 0$ everywhere: $U^V$ is strictly convex on $[0,1]$ and attains its maximum at a corner.
 
The first derivative is $L(\beta_h) = \frac{1}{2}WW' + \frac{Q}{t}y_A'$. Substituting:
\[
L = \frac{\delta}{t}\!\left[\frac{W/2}{m_H^2} - \frac{\alpha W/2 + Q}{m_A^2}\right] = \frac{\delta}{t}\!\left[\frac{r^*}{m_H^2} - \frac{q_A^*}{m_A^2}\right].
\]
For $Q = 0$: $L = \frac{\delta r^*}{t}\!\left(\frac{1}{m_H^2} - \frac{\alpha}{m_A^2}\right)$, which is positive for typical parameters (since $\alpha \leq 1$), pushing toward $\beta_h = 1$. For $Q$ large: the term $-Q\delta/(tm_A^2)$ dominates, making $L < 0$ and pushing toward $\beta_h = 0$. By continuity, there exists a threshold $Q^*$ (defined implicitly by $U^V(0) = U^V(1)$) separating the two regimes.

 
\medskip\noindent\textbf{Part~(ii): $k > \bar\Gamma  := \max_{\beta_h \in [0,1]} \Gamma(\beta_h)$.}
The penalised second derivative $V''^V = \Gamma - k < 0$ everywhere, so $U^V$ is strictly concave with a unique maximum. The first derivative is $dU^V/d\beta_h = L(\beta_h) - k(\beta_h - \frac{1}{2})$. At the boundaries:
\[
\left.\frac{dU^V}{d\beta_h}\right|_0 = L(0) + \frac{k}{2}, \qquad \left.\frac{dU^V}{d\beta_h}\right|_1 = L(1) - \frac{k}{2}.
\]
 
\emph{Sub-case~(ii-a): $k > k^* = \max\{\bar\Gamma, -2L(0), 2L(1)\}$.}
Then $L(0) + k/2 > 0$ and $L(1) - k/2 < 0$: the function increases at the left boundary and decreases at the right. By strict concavity, the unique maximum is interior, satisfying $L(\beta_h^*) = k(\beta_h^* - \frac{1}{2})$.
 
\emph{Sub-case~(ii-b): $L(0) \leq -k/2$.}
$dU^V/d\beta_h|_0 \leq 0$. Strict concavity implies $U^V$ is decreasing on $[0,1]$, so $\beta_h^* = 0$.
 
\emph{Sub-case~(ii-c): $L(1) \geq k/2$.}
$dU^V/d\beta_h|_1 \geq 0$. Strict concavity implies $U^V$ is increasing on $[0,1]$, so $\beta_h^* = 1$.
 
\medskip\noindent\textbf{Part~(iii): $0 < k \leq \bar\Gamma$.}
$V''^V = \Gamma - k$ may change sign. The utility function is continuous on $[0,1]$, so the Extreme Value Theorem guarantees a global maximum. It is found by solving $L(\beta_h) = k(\beta_h - \frac{1}{2})$ for all stationary points, evaluating $U^V$ at each, comparing with $U^V(0)$ and $U^V(1)$, and selecting the largest. \hfill$\square$
 
\subsection*{Proof of Proposition~\ref{prop:rstar-eng}}
 
Substitute $q_h^* = r$ and $q_A^* = \alpha r + Q$ into the engagement payoff (omitting the penalty):
\[
u_p^E = \frac{(\alpha r + Q)^2}{tm_A} + \frac{r^2}{tm_H} - r^2 = \frac{\alpha^2 r^2 + 2\alpha Qr + Q^2}{tm_A} + r^2\!\left(\frac{1}{tm_H} - 1\right).
\]
Collecting powers of $r$:
\[
u_p^E = r^2\!\left(\frac{\alpha^2}{tm_A} + \frac{1}{tm_H} - 1\right) + \frac{2\alpha Qr}{tm_A} + \frac{Q^2}{tm_A}.
\]
The coefficient of $r^2$ has the sign of $\alpha^2/(tm_A) + 1/(tm_H) - 1$, or equivalently, of $\alpha^2 m_H + m_A - tm_Am_H = -\Delta$.
 
\medskip\noindent\textbf{Part~(a): $\Delta > 0$.}
The $r^2$ coefficient is negative, confirming strict concavity. The FOC $\partial u_p^E/\partial r = 0$ gives:
\[
2r\!\left(\frac{\alpha^2}{tm_A} + \frac{1}{tm_H} - 1\right) + \frac{2\alpha Q}{tm_A} = 0.
\]
Multiplying by $tm_Am_H/2$ and solving:
\[
-r\,\Delta + \alpha Qm_H = 0 \quad\Longrightarrow\quad r^* = \frac{\alpha Qm_H}{\Delta}.
\]
Since $\Delta > 0$, $\alpha > 0$, and $m_H > 0$: $r^* > 0$ iff $Q > 0$.
 
For the sufficient condition: $\Delta = tm_Am_H - \alpha^2 m_H - m_A$. At $\beta_h = 1$ (where $m_A = 1$, $m_H = 1-\delta$): $\Delta = t(1-\delta) - \alpha^2(1-\delta) - 1 = (1-\delta)(t-\alpha^2) - 1$. This is positive iff $t > \alpha^2 + 1/(1-\delta)$.
 
\medskip\noindent\textbf{Part~(b): $\Delta \leq 0$.}
The $r^2$ coefficient is non-negative, so $u_p^E$ is convex (or linear) in $r$. The FOC gives a minimum, and the platform would set $r \to \infty$ (unbounded subsidization), so the model requires $\Delta > 0$. \hfill$\square$
 
\subsection*{Proof of Proposition~\ref{prop:beta-eng}}
 
\noindent\textbf{Reduction to a quadratic in $r$.}
Substituting the Stage-2 equilibrium $q_h^* = r$ and $q_A^* = \alpha r + Q$ into the demand functions and engagement payoff (omitting the $r$-independent penalty):
\[
u_p^E(r) = q_A^*\,D_A + q_h^*\,D_h - r\,q_h^* = \frac{(\alpha r + Q)^2}{tm_A} + \frac{r^2}{tm_H} - r^2.
\]
Expanding $(\alpha r + Q)^2 = \alpha^2 r^2 + 2\alpha Qr + Q^2$ and grouping by powers of~$r$:
\[
u_p^E(r) = r^2\!\left[\frac{\alpha^2}{tm_A} + \frac{1}{tm_H} - 1\right] + r\cdot\frac{2\alpha Q}{tm_A} + \frac{Q^2}{tm_A}.
\]
Multiplying the $r^2$ coefficient by $tm_Am_H > 0$:
\[
tm_Am_H\!\left[\frac{\alpha^2}{tm_A} + \frac{1}{tm_H} - 1\right] = \alpha^2 m_H + m_A - tm_Am_H = -\Delta.
\]
Hence the coefficient of $r^2$ equals $-\Delta/(tm_Am_H)$, and
\[
u_p^E(r) = \underbrace{\frac{-\Delta}{tm_Am_H}}_{a\,<\,0}\,r^2 + \underbrace{\frac{2\alpha Q}{tm_A}}_{b}\,r + \underbrace{\frac{Q^2}{tm_A}}_{c},
\]
with $a < 0$ by the assumption $\Delta > 0$. This is a concave quadratic in $r$, with unique maximiser
\[
r^* = -\frac{b}{2a} = \frac{2\alpha Q/(tm_A)}{2\Delta/(tm_Am_H)} = \frac{\alpha\,Q\,m_H}{\Delta}.
\]
 
\medskip\noindent\textbf{Utility function.}
Using the formula $U = c - b^2/(4a)$ for the maximum of $ar^2 + br + c$ with $a < 0$:
\[
-\frac{b^2}{4a} = \frac{\bigl(2\alpha Q/(tm_A)\bigr)^2}{4\,\Delta/(tm_Am_H)} = \frac{4\alpha^2 Q^2\,tm_Am_H}{4\Delta\,(tm_A)^2} = \frac{\alpha^2 Q^2 m_H}{tm_A\,\Delta}.
\]
Adding $c = Q^2/(tm_A)$:
\[
U^E\big|_{k=0} = \frac{Q^2}{tm_A} + \frac{\alpha^2 Q^2 m_H}{tm_A\,\Delta} = \frac{Q^2}{tm_A}\cdot\frac{\Delta + \alpha^2 m_H}{\Delta}.
\]
Using the identity
\[
\Delta + \alpha^2 m_H = (tm_Am_H - \alpha^2 m_H - m_A) + \alpha^2 m_H = tm_Am_H - m_A = m_A(tm_H - 1):
\]
\[
U^E\big|_{k=0} = \frac{Q^2\cdot m_A(tm_H-1)}{tm_A\,\Delta} = \frac{Q^2(tm_H-1)}{t\,\Delta}.
\]
Adding the penalty gives~\eqref{eq:VE}.
 
\medskip\noindent\textbf{Marginal promotion effect.}
By the Envelope Theorem, since $\partial u_p^E/\partial r = 0$ at $r^*$, the total derivative with respect to $\beta_h$ reduces to the partial derivative at fixed $r = r^*$:
\[
\frac{dU^E}{d\beta_h} = \left.\frac{\partial u_p^E}{\partial\beta_h}\right|_{r=r^*} - k\!\left(\beta_h-\frac{1}{2}\right).
\]
Holding $r$ fixed, $q_A = \alpha r + Q$ does not depend on $\beta_h$, so in $u_p^E = q_A^2/(tm_A) + r^2/(tm_H) - r^2$ the dependence on $\beta_h$ enters only through $m_A$ and $m_H$. Using $m_A' = \delta$ and $m_H' = -\delta$:
\[
\frac{\partial}{\partial\beta_h}\!\left(\frac{q_A^2}{tm_A}\right) = q_A^2\cdot\frac{-m_A'}{tm_A^2} = -\frac{q_A^2\,\delta}{tm_A^2},
\]
\[
\frac{\partial}{\partial\beta_h}\!\left(\frac{r^2}{tm_H}\right) = r^2\cdot\frac{-m_H'}{tm_H^2} = \frac{r^2\,\delta}{tm_H^2}.
\]
The last term $-r^2$ is independent of $\beta_h$. Summing and evaluating at $r = r^*$:
\[
L(\beta_h) := \left.\frac{\partial u_p^E}{\partial\beta_h}\right|_{r=r^*} = \frac{\delta}{t}\!\left[\frac{(r^*)^2}{m_H^2} - \frac{(q_A^*)^2}{m_A^2}\right].
\]


From Step~1, $r^* = \alpha Qm_H/\Delta$. For $q_A^*$:
\[
q_A^* = \alpha r^* + Q = \frac{\alpha^2 Qm_H}{\Delta} + Q = \frac{Q(\alpha^2 m_H + \Delta)}{\Delta} = \frac{Q\,m_A(tm_H-1)}{\Delta},
\]

and substituting into the expression for $L$:
\[
L(\beta_h) = \frac{\delta}{t}\!\left[\frac{\alpha^2 Q^2}{\Delta^2} - \frac{Q^2(tm_H-1)^2}{\Delta^2}\right] = \frac{\delta\,Q^2}{t\,\Delta^2}\!\left[\alpha^2 - (tm_H-1)^2\right].
\]
 
\medskip\noindent\textbf{Sign of $L$.}
Factoring: $\alpha^2 - (tm_H-1)^2 = [\alpha + (tm_H-1)][\alpha - (tm_H-1)]$. Since $\Delta > 0$ implies $tm_H > 1$, the first factor is always positive. Hence $\operatorname{sign}(L) = \operatorname{sign}[(1+\alpha) - tm_H]$.
 
Since $tm_H = t(1-\beta_h\delta)$ is strictly decreasing, $L > 0 \Leftrightarrow \beta_h > \beta_c$ where $\beta_c = (t-1-\alpha)/(t\delta)$. Three regimes arise:
\begin{itemize}
\item $t \leq 1+\alpha$: $\beta_c \leq 0$, so $L \geq 0$ on $[0,1]$.
\item $t \geq (1+\alpha)/(1-\delta)$: $\beta_c \geq 1$, so $L \leq 0$ on $[0,1]$.
\item $1+\alpha < t < (1+\alpha)/(1-\delta)$: $\beta_c \in (0,1)$, with $L < 0$ on $[0,\beta_c)$ and $L > 0$ on $(\beta_c,1]$.
\end{itemize}
 
\medskip\noindent\textbf{Part~(i): $k = 0$.}
$dU^E/d\beta_h = L(\beta_h)$.

 
\emph{Sub-case~(i-a): $t \leq 1+\alpha$.}
$L \geq 0$ on $[0,1]$: $U^E$ is non-decreasing, so $\beta_h^* = 1$.
 
\emph{Sub-case~(i-b): $t \geq (1+\alpha)/(1-\delta)$.}
$L \leq 0$ on $[0,1]$: $U^E$ is non-increasing, so $\beta_h^* = 0$.
 
\emph{Sub-case~(i-c): $1+\alpha < t < (1+\alpha)/(1-\delta)$.}
$L < 0$ on $[0,\beta_c)$ and $L > 0$ on $(\beta_c,1]$: $U^E$ decreases then increases, with a minimum at $\beta_c$. The maximum is at a corner, $\beta_h^* \in \{0,1\}$, chosen by comparing $U^E(0)$ and $U^E(1)$. \hfill$\square$

 
\medskip\noindent\textbf{Part~(ii): $k > 0$.}
The derivative is $G(\beta_h) = L(\beta_h) - k(\beta_h - \tfrac{1}{2})$, with boundary values
\[
G(0) = L(0) + \frac{k}{2}, \qquad G(1) = L(1) - \frac{k}{2}.
\]

 
%%-------------------------------------------------------------
\medskip\noindent\textbf{Regime~A: $t \leq 1 + \alpha$ (so $L \geq 0$ everywhere).}
 
Since $L(\beta_h) \geq 0$ for all $\beta_h$, the boundary values satisfy $G(0) = L(0) + k/2 > 0$ always (both terms positive). The behavior hinges on the right boundary:
\begin{itemize}
\item If $k > 2L(1)$: $G(1) = L(1) - k/2 < 0$. Since $G(0) > 0$ and $G(1) < 0$, the IVT yields an interior zero. At this zero, $F(\beta_h^*) = k(\beta_h^* - \frac{1}{2}) > 0$ (since $L \geq 0$), which forces $\beta_h^* > \frac{1}{2}$: the platform tilts toward human content but the penalty prevents full promotion.
 
\item If $k \leq 2L(1)$: $G(1) \geq 0$. We claim $G(\beta_h) \geq 0$ on $[0,1]$, so $U^E$ is non-decreasing and $\beta_h^* = 1$. To see this: for $\beta_h \leq \frac{1}{2}$, $L(\beta_h) \geq 0$ and $-k(\beta_h - \frac{1}{2}) \geq 0$, so $G \geq 0$. For $\beta_h > \frac{1}{2}$: in this regime $L$ is non-negative and increasing in $\beta_h$ . The penalty $-k(\beta_h - \frac{1}{2})$ is linear and its steepest value is $-k/2$ at $\beta_h = 1$. Since $G(1) = L(1) - k/2 \geq 0$, $l$ dominates the penalty throughout $[\frac{1}{2}, 1]$, so $G \geq 0$.
\end{itemize}
 
Summary: the threshold is $k_A = 2L(1)$. For $k \leq k_A$, $\beta_h^* = 1$. For $k > k_A$, $\beta_h^* \in (\frac{1}{2}, 1)$.
 
%%-------------------------------------------------------------
\medskip\noindent\textbf{Regime~B: $t \geq (1+\alpha)/(1-\delta)$ (so $L \leq 0$ everywhere).}
 
Since $L(\beta_h) \leq 0$ for all $\beta_h$, $G(1) = L(1) - k/2 < 0$ always (both terms negative). The behavior hinges on the left boundary:
\begin{itemize}
\item If $k > -2L(0)$: $G(0) = L(0) + k/2 > 0$. Since $G(0) > 0$ and $G(1) < 0$, the IVT yields an interior zero. At this zero, $L(\beta_h^*) = k(\beta_h^* - \frac{1}{2}) < 0$, which forces $\beta_h^* < \frac{1}{2}$.
 
\item If $k \leq -2L(0)$: $G(0) \leq 0$. We claim $G(\beta_h) \leq 0$ on $[0,1]$, so $U^E$ is non-increasing and $\beta_h^* = 0$. For $\beta_h \geq \frac{1}{2}$: $L(\beta_h) \leq 0$ and $-k(\beta_h - \frac{1}{2}) \leq 0$, so $G \leq 0$. For $\beta_h < \frac{1}{2}$: in this regime $L$ is non-positive and becomes more negative as $\beta_h$ decreases. The penalty $-k(\beta_h - \frac{1}{2})$ is positive but bounded by $k/2$. Since $G(0) = L(0) + k/2 \leq 0$, the negative $L$ dominates the penalty throughout $[0, \frac{1}{2})$, so $G \leq 0$.
\end{itemize}
 
Summary: the threshold is $k_B = 2|L(0)|$. For $k \leq k_B$, $\beta_h^* = 0$. For $k > k_B$, $\beta_h^* \in (0, \frac{1}{2})$.
 
%%-------------------------------------------------------------
\medskip\noindent\textbf{Regime~C: $1+\alpha < t < (1+\alpha)/(1-\delta)$ (so $L$ changes sign at $\beta_c$).}
 
Here $L(0) < 0$ and $L(1) > 0$, so $G(0) = L(0) + k/2$ and $G(1) = L(1) - k/2$ can each be positive or negative. Define $k^{\star} = \max\{2|L(0)|,\; 2L(1)\}$.
 
\begin{itemize}
\item If $k > k^{\star}$: $G(0) > 0$ and $G(1) < 0$. The IVT yields an interior $\beta_h^* \in (0,1)$.
 
\item If $k \leq 2|L(0)|$ and $k \leq 2L(1)$, i.e., $k \leq \min\{2|L(0)|, 2L(1)\}$: $G(0) \leq 0$ and $G(1) \geq 0$. The utility is non-increasing at the left and non-decreasing at the right, reproducing the U-shape from Part~(i-c). The global maximum is either at a corner $\beta_h^* \in \{0,1\}$ or at an interior stationary point, determined by comparing $U^E(0)$ and $U^E(1)$.
 
\item If $k \leq 2|L(0)|$ and $k > 2L(1)$: $G(0) \leq 0$ and $G(1) < 0$. The utility is non-increasing at both boundaries. The global optimum is $\beta_h^* = 0$ or the interior peak, determined by comparing $U^E$ values.
 
\item If $k > 2|L(0)|$ and $k \leq 2L(1)$: $G(0) > 0$ and $G(1) \geq 0$. The utility is non-decreasing at both boundaries. The global optimum is $\beta_h^* = 1$ or an interior peak, determined by comparing $U^E$ values.
\end{itemize}
 
%%-------------------------------------------------------------
\medskip\noindent\textbf{Convergence as $k \to \infty$.}
Let $M = \max_{\beta_h}|L(\beta_h)| < \infty$. For $k > 2M$, we have $G(0) > 0$ and $G(1) < 0$ regardless of $t$-regime, so the interior case applies. Any solution satisfies $|\beta_h^* - \frac{1}{2}| \leq M/k \to 0$. \hfill$\square$
 
 


\vfill
\pagebreak

 
 

\end{document}